\section{Related Work}

\subsection{Human--AI interaction and intelligible control}

Human--AI interaction research repeatedly shows that a capable model is not
enough: an interface must communicate what the system can do, expose status and
uncertainty, support correction, and help people recover from errors
\cite{amershi2019guidelines,yang2020reexamining,weisz2024design}. Explanation
needs depend on the questions people ask about a system
\cite{liao2020questioning}, while explanations do not automatically improve
joint performance \cite{bansal2021whole}. These findings motivate three design
choices in our system: agent responsibility is inspectable, routing is reported
with a reason, and unresolved spatial references fail visibly instead of being
silently guessed. We do not propose new general HAI guidelines or infer
subjective benefit from these mechanisms; we test whether the implementation
obeys its declared interaction contract.

\subsection{Embodied agents and spatial reference}

Embodied conversational agents established how language, gaze, gesture, and a
visible body combine in an interface \cite{cassell1999embodiment}. Museum
guides further demonstrate the practical and educational importance of agent
behavior in a shared place \cite{kopp2005museum,bickmore2011museum}. Recent
generative agents add memory, reflection, and social behavior
\cite{park2023generative}, while multi-agent frameworks coordinate
language-model roles through conversation \cite{li2023camel}. Those systems
motivate flexible agent behavior, but they do not by themselves specify which
Unity object an utterance concerns or provide a requirements-to-runtime
validation trace.

Situated-language research directly studies this reference problem. Approaches
combine recognition of a user's plan and referring expression
\cite{gorniak2005grounding}, connect directions to perceived environments
\cite{kollar2010directions}, maintain symbolic anchors for human--robot
discourse \cite{lemaignan2012grounding}, or learn perceptually grounded word
meanings for incremental reference resolution
\cite{kennington2015reference}. \systemname{} does not introduce another
general grounding model. It supplies a system boundary around a chosen
selection or grounding mechanism: candidates are bound to stable model
entities, ambiguity is explicit, responsibility is constrained, and the
resolved entity remains available to runtime assertions.

\subsection{Conversational authoring and agents in XR}

Several recent systems make language models useful inside interactive 3D
environments. LLMR prompts and edits interactive mixed-reality worlds in real
time and includes planning, scene understanding, and self-debugging
\cite{delatorre2024llmr}. Tell-XR lets end users author event--condition--action
automations conversationally in AR and VR
\cite{carcangiu2025tellxr}. CUIfy packages the integration of LLM-powered
conversational agents in XR \cite{buldu2025cuify}. AgentHands studies
interactive gestures for spatially grounded XR conversations
\cite{liu2026agenthands}.

These systems rule out a claim that combining Unity or XR with LLM agents is
itself novel. Our focus is complementary. We make the generated interaction
architecture inspectable across spatial entities, agent responsibilities,
capabilities, backend actions, and validation evidence. Compared with LLMR and
Tell-XR, the target is not free-form world generation or individual automation
rules, but lifecycle-wide consistency of an embodied multi-agent interface.
Compared with CUIfy and AgentHands, the contribution is an authoring and
verification layer rather than an integration library or gesture-generation
method.

\subsection{Model-based interface engineering and traceability}

Model-based interface development has long addressed the mapping from abstract
tasks and domain concepts to concrete interfaces
\cite{puerta1999framework}. The Cameleon framework organizes transformations
across abstraction levels and deployment targets
\cite{calvary2003cameleon}; UsiXML supports multiple development paths
\cite{limbourg2005usixml}; MARIA describes service-oriented interfaces at
multiple abstraction levels \cite{paterno2009maria}; and OpenUIDL supports
runtime, omni-channel interfaces \cite{moldovan2020openuidl}. FunctionalMLDS is
narrower: its domain is spatial, agentic interaction, and its platform mappings
connect Unity entities and events to agent capabilities and backend operations.

The design also builds on requirements traceability: a trace is useful when it
helps locate and reason about consequences across artifacts
\cite{gotel1994traceability}, and empirical work shows that developers can
benefit from trace links during maintenance and evolution
\cite{maeder2015benefit}. We therefore evaluate concrete change and fault tasks
rather than treating the presence of trace links as evidence of benefit.

\subsection{Positioning}

Prior work separately addresses model-based interface generation,
conversational XR authoring, spatial grounding, embodied agents, and LLM-agent
orchestration. \systemname{} connects these concerns through one executable
interaction contract linking a user scenario, selected spatial entity,
responsible agent, capability, Unity/backend action, and observed assertion.
The empirical contribution is correspondingly bounded: we test reference and
routing correctness, change consistency, fault detection, and trace
completeness against direct artifact wiring. We do not claim the first XR agent
system, a new grounding algorithm, or a general-purpose user-interface
description language.
